\chapter{An Elliptic Integral and Its Differential Equation}
\label{ch:elliptic}

A final example ties together rational parametrization, algebraic-value
computations, integration, series, and a variable-coefficient fundamental
matrix.  The point is not to develop all elliptic functions, but to give
several equivalent ways of computing one useful family.

\section{Remove an endpoint singularity before computing}
For rational $0\le m<1$, the familiar expression is
\[
 K(m)=\int_0^1\frac{dx}{\sqrt{(1-x^2)(1-mx^2)}}.
\]
We do not begin by assuming this improper integral exists.  Use the rational
circle parameter $x=2u/(1+u^2)$.  For $0\le u<1$,
\[
 \sqrt{1-x^2}=\frac{1-u^2}{1+u^2},\qquad
 \frac{dx}{du}=\frac{2(1-u^2)}{(1+u^2)^2}.
\]
Cancellation suggests the bounded integral
\[
 \boxed{K(m)=\int_0^1
       \frac{2\,du}{\sqrt{(1+u^2)^2-4mu^2}}.}
\]
This is our first definition.  On $0\le m\le r<1$ its squared denominator
is at least $(1-r)(1+u^2)^2\ge1-r$.  Positive-root bisection and separated
reciprocals therefore give value computations, absolute bounds, and
continuity moduli for the integrand.

For truncated $u$-intervals the substitution theorem proves agreement with
the corresponding truncated original integral.  The omitted transformed
interval of length $\varepsilon$ contributes at most
$2\varepsilon/\sqrt{1-r}$, with any rational upper bound for this constant.
It follows that the original improper expression is an equivalent
computation, now accompanied by its missing tail estimate.

At $m=0$ the integrand is $2/(1+u^2)$, so
$K(0)=2A(1)=\pi/2$.

\section{A power series from the same integral}
Put $c_n=\binom{2n}{n}/4^n$, as in
Chapter~\ref{ch:infinite-series}.  The transformed integrand is
\[
 \frac2{1+u^2}\sum_{n=0}^{\infty}
       c_n m^n\left(\frac{2u}{1+u^2}\right)^{2n}.
\]
Since $0\le 2u/(1+u^2)\le1$ and $c_n\le1$, truncating after $N$ has
uniform error at most $2r^{N+1}/(1-r)$.  Integrating finite prefixes and
using that error bound is therefore justified.

Let
\[
 I_n=\int_0^1\frac2{1+u^2}
          \left(\frac{2u}{1+u^2}\right)^{2n}du.
\]
The angle substitution $x=2A(u)/\pi$ gives
$I_n=\pi\int_0^{1/2}S(x)^{2n}dx$.  Differentiate the finite product
$S^{2n-1}C$ using the already proved trigonometric derivatives and product
rule.  Its endpoint values vanish for $n\ge1$, so integration by parts
gives
\[
 2n I_n=(2n-1)I_{n-1},\qquad I_0=\pi/2.
\]
The powers and products have uniform finite increment estimates obtained
by multiplying those of $S$ and $C$, so the telescoping identity used here
is certified on the whole interval.  Hence $I_n=(\pi/2)c_n$ and
\[
 \boxed{K(m)=\frac\pi2\sum_{n=0}^{\infty}c_n^2m^n.}
\]
A tail bound is $(\pi/2)r^{N+1}/(1-r)$, which can be replaced by the
rational bound $2r^{N+1}/(1-r)$.

For example, at $m=1/2$ the first five terms of the normalized quantity
$2K(m)/\pi$ are
\[
 1+\frac18+\frac9{256}+\frac{25}{2048}
          +\frac{1225}{262144}.
\]
The remaining normalized tail is at most $1/16$.  More precise boxes
are obtained by the same rational recurrence, not by a new existence
argument.  Rectangle quadrature and this series now compute the same number.

\section{The differential equation is a coefficient identity}
Writing $a_n=(\pi/2)c_n^2$, the recurrence is
\[
 (n+1)^2a_{n+1}=(n+1/2)^2a_n.
\]
The differentiated tails converge uniformly on $m\le r<1$, by the power
series estimates.  Comparing coefficients therefore gives
\[
 \boxed{m(1-m)K''+(1-2m)K'-\tfrac14K=0,}
\]
with $K(0)=\pi/2$ and $K'(0)=\pi/8$.
This is proved from finite recurrences and their tails, not from an abstract
classification of period functions.

On an interval strictly inside $(0,1)$, put $Y=(K,K')^{\mathsf T}$.  It
satisfies the linear system
\[
 Y'=\begin{pmatrix}
 0&1\\
 \dfrac1{4m(1-m)}&\dfrac{2m-1}{m(1-m)}
 \end{pmatrix}Y.
\]
For example, on $[1/4,3/4]$ the row-sum norm of this coefficient matrix is
at most $4$.  The fundamental-matrix construction therefore has the
factorial majorant of the preceding chapter.  The series supplies
$Y(1/2)$, and propagation supplies $Y(3/4)$; integral-equation uniqueness
identifies it with direct series evaluation there.

The coefficient matrix is singular at $m=0$ and $m=1$.  We have not
applied its bounded-coefficient theorem across either point.  The series
handles the specified regular solution at $m=0$; each evolution computation
uses an explicitly separated interval.

\section{What the example connects}
We now have three actual computations of $K$: bounded quadrature of an
algebraic integrand, a rational coefficient series times $\pi/2$, and
fundamental-matrix propagation from a computed initial state.  Their
comparisons are proofs: rational substitution with an endpoint tail,
termwise finite integration with a uniform remainder, and uniqueness for
an integral equation.

The same matrix also computes a response to any finite list of impulses
in its two state coordinates on a separated interval.  Thus delta inputs
need no new theory for this example.  They are finite state updates
propagated by a computation already constructed for a concrete special
function.  This is the recurring pattern of the book: retain useful
representations, prove their equivalence, and reuse the resulting
computations in the next problem.
